% ==============================================================================
% Formation C# & SQL — Module 001 : Les Bases
% Fichier : src/P1_05_a.tex
% Sujet   : Chapitre 5 : Les Boucles
% ==============================================================================

\newpage
\section{Les Boucles (Structures Itératives)}
\marginpar{\tiny\color{gray}P1\_05\_a.tex}

Les boucles, ou structures itératives, sont des mécanismes de contrôle de flux permettant d'exécuter un bloc d'instructions de manière répétée tant qu'une condition booléenne demeure évaluée à \texttt{true}. Elles sont indispensables en ingénierie logicielle pour l'automatisation de tâches répétitives, le parcours de structures de données ou l'attente d'événements asynchrones.

\subsection{La boucle \texttt{for}}
La boucle \texttt{for} est privilégiée lorsque le nombre d'itérations est connu à l'avance ou déterminé par une borne fixe. Elle encapsule l'initialisation de l'index, la condition de maintien et l'incrémentation (ou décrémentation) dans une seule ligne de déclaration.

\begin{lstlisting}[style=csharpstyle]
for (int i = 0; i < 5; i++)
{
    Console.WriteLine($"Itération numéro : {i}");
}
\end{lstlisting}

\textbf{Analyse de l'exécution :}
\begin{itemize}
    \item \textbf{Initialisation} (\texttt{int i = 0}) : Exécutée une seule fois au démarrage de la boucle.
    \item \textbf{Condition} (\texttt{i < 5}) : Évaluée avant chaque itération. Si elle est \texttt{false}, la boucle s'arrête et le programme continue à l'instruction suivante.
    \item \textbf{Incrémentation} (\texttt{i++}) : Exécutée à la fin de chaque itération, juste avant de réévaluer la condition.
\end{itemize}

\subsection{La boucle \texttt{while}}
La boucle \texttt{while} évalue sa condition de maintien \textit{avant} d'exécuter son bloc d'instructions. Elle est typiquement employée lorsque le nombre d'itérations est indéterminé, dépendant d'un état dynamique du système ou de la lecture d'un flux de données (ex: lecture de fichiers, écoute réseau).

\begin{lstlisting}[style=csharpstyle]
int bufferSize = 0;
while (bufferSize < 5)
{
    Console.WriteLine($"Traitement du bloc mémoire {bufferSize}...");
    bufferSize++;
}
\end{lstlisting}
Contrairement au \texttt{for}, l'initialisation de la variable de contrôle est externe à la boucle, et son altération doit obligatoirement figurer à l'intérieur du bloc pour éviter une boucle infinie.

\subsection{La boucle \texttt{do...while}}
Contrairement à la boucle \texttt{while}, la variante \texttt{do...while} évalue la condition \textit{après} l'exécution du bloc. Cela garantit que les instructions internes seront exécutées au strict minimum une fois, indépendamment de l'état initial.

\begin{lstlisting}[style=csharpstyle]
int retryCount = 0;
do
{
    Console.WriteLine($"Tentative de connexion réseau (Essai {retryCount + 1})");
    retryCount++;
} while (retryCount < 5);
\end{lstlisting}
Cette boucle est très utile pour les scénarios nécessitant au moins une action initiale, tels que les tentatives de connexion (\textit{retry policy}) ou l'exécution d'un menu interactif.

\subsection{La boucle \texttt{foreach}}
La structure \texttt{foreach} offre une syntaxe abstraite et sécurisée pour itérer sur l'ensemble des éléments d'une collection (tableaux, listes) implémentant l'interface \texttt{IEnumerable}. Elle élimine la gestion manuelle de l'index et prévient les erreurs de dépassement.

\begin{lstlisting}[style=csharpstyle]
int[] tcpPorts = { 80, 443, 8080, 1433 };
foreach (int port in tcpPorts)
{
    Console.WriteLine($"Analyse du port : {port}");
}
\end{lstlisting}
Bien que cette boucle produise un code plus concis et robuste, la variable d'itération (\texttt{port} dans l'exemple) est en lecture seule. Vous ne pouvez pas modifier directement l'élément du tableau via cette variable.

\subsection{Contrôle de Flux : \texttt{break} et \texttt{continue}}
Au sein de toute structure itérative, il est possible d'altérer le flux standard à l'aide de deux mots-clés :
\begin{itemize}
    \item \texttt{break} : Interrompt immédiatement et définitivement la boucle englobante, forçant le programme à passer à l'instruction qui suit la boucle.
    \item \texttt{continue} : Interrompt l'itération courante et passe directement à l'évaluation de la condition pour démarrer l'itération suivante.
\end{itemize}

\begin{tcolorbox}[colback=backcolour,colframe=keywordcolor,title=\textbf{Pièges Courants \& Erreurs}]
\begin{itemize}
    \item \textbf{La boucle infinie} : Oublier d'incrémenter la variable d'arrêt ou de modifier la condition dans une boucle \texttt{while}, entraînant le blocage complet du thread (\textit{Thread Starvation}).
    \item \textbf{Modification de la collection} : Tenter d'ajouter ou de supprimer un élément d'une liste qui est actuellement en train d'être parcourue par un \texttt{foreach} déclenche immanquablement une \texttt{InvalidOperationException} à l'exécution.
    \item \textbf{Erreur de borne (\textit{Off-by-one error})} : Écrire \texttt{for (int i = 0; i <= notes.Length; i++)} au lieu de \texttt{i < notes.Length}. Puisque les éléments d'un tableau de taille $N$ sont indexés de $0$ à $N-1$, accéder à l'index $N$ (\texttt{notes[notes.Length]}) lève systématiquement une \texttt{IndexOutOfRangeException} fatale.
\end{itemize}
\end{tcolorbox}

\subsection{Synthèse}
\begin{itemize}
    \item \textbf{\texttt{for}} : À privilégier pour les itérations indexées dont le nombre est prévisible et fini.
    \item \textbf{\texttt{while}} : Idéal pour les processus en attente d'un état ou lisant un flux de données continu.
    \item \textbf{\texttt{do...while}} : Utile lorsqu'un premier traitement est impératif avant toute vérification de condition.
    \item \textbf{\texttt{foreach}} : La méthode la plus élégante et robuste pour parcourir des structures de données en toute sécurité.
\end{itemize}
